% Options for packages loaded elsewhere
\PassOptionsToPackage{unicode}{hyperref}
\PassOptionsToPackage{hyphens}{url}
\documentclass[
]{article}
\usepackage{xcolor}
\usepackage[margin=1in]{geometry}
\usepackage{amsmath,amssymb}
\setcounter{secnumdepth}{-\maxdimen} % remove section numbering
\usepackage{iftex}
\ifPDFTeX
  \usepackage[T1]{fontenc}
  \usepackage[utf8]{inputenc}
  \usepackage{textcomp} % provide euro and other symbols
\else % if luatex or xetex
  \usepackage{unicode-math} % this also loads fontspec
  \defaultfontfeatures{Scale=MatchLowercase}
  \defaultfontfeatures[\rmfamily]{Ligatures=TeX,Scale=1}
\fi
\usepackage{lmodern}
\ifPDFTeX\else
  % xetex/luatex font selection
\fi
% Use upquote if available, for straight quotes in verbatim environments
\IfFileExists{upquote.sty}{\usepackage{upquote}}{}
\IfFileExists{microtype.sty}{% use microtype if available
  \usepackage[]{microtype}
  \UseMicrotypeSet[protrusion]{basicmath} % disable protrusion for tt fonts
}{}
\makeatletter
\@ifundefined{KOMAClassName}{% if non-KOMA class
  \IfFileExists{parskip.sty}{%
    \usepackage{parskip}
  }{% else
    \setlength{\parindent}{0pt}
    \setlength{\parskip}{6pt plus 2pt minus 1pt}}
}{% if KOMA class
  \KOMAoptions{parskip=half}}
\makeatother
\usepackage{color}
\usepackage{fancyvrb}
\newcommand{\VerbBar}{|}
\newcommand{\VERB}{\Verb[commandchars=\\\{\}]}
\DefineVerbatimEnvironment{Highlighting}{Verbatim}{commandchars=\\\{\}}
% Add ',fontsize=\small' for more characters per line
\usepackage{framed}
\definecolor{shadecolor}{RGB}{248,248,248}
\newenvironment{Shaded}{\begin{snugshade}}{\end{snugshade}}
\newcommand{\AlertTok}[1]{\textcolor[rgb]{0.94,0.16,0.16}{#1}}
\newcommand{\AnnotationTok}[1]{\textcolor[rgb]{0.56,0.35,0.01}{\textbf{\textit{#1}}}}
\newcommand{\AttributeTok}[1]{\textcolor[rgb]{0.13,0.29,0.53}{#1}}
\newcommand{\BaseNTok}[1]{\textcolor[rgb]{0.00,0.00,0.81}{#1}}
\newcommand{\BuiltInTok}[1]{#1}
\newcommand{\CharTok}[1]{\textcolor[rgb]{0.31,0.60,0.02}{#1}}
\newcommand{\CommentTok}[1]{\textcolor[rgb]{0.56,0.35,0.01}{\textit{#1}}}
\newcommand{\CommentVarTok}[1]{\textcolor[rgb]{0.56,0.35,0.01}{\textbf{\textit{#1}}}}
\newcommand{\ConstantTok}[1]{\textcolor[rgb]{0.56,0.35,0.01}{#1}}
\newcommand{\ControlFlowTok}[1]{\textcolor[rgb]{0.13,0.29,0.53}{\textbf{#1}}}
\newcommand{\DataTypeTok}[1]{\textcolor[rgb]{0.13,0.29,0.53}{#1}}
\newcommand{\DecValTok}[1]{\textcolor[rgb]{0.00,0.00,0.81}{#1}}
\newcommand{\DocumentationTok}[1]{\textcolor[rgb]{0.56,0.35,0.01}{\textbf{\textit{#1}}}}
\newcommand{\ErrorTok}[1]{\textcolor[rgb]{0.64,0.00,0.00}{\textbf{#1}}}
\newcommand{\ExtensionTok}[1]{#1}
\newcommand{\FloatTok}[1]{\textcolor[rgb]{0.00,0.00,0.81}{#1}}
\newcommand{\FunctionTok}[1]{\textcolor[rgb]{0.13,0.29,0.53}{\textbf{#1}}}
\newcommand{\ImportTok}[1]{#1}
\newcommand{\InformationTok}[1]{\textcolor[rgb]{0.56,0.35,0.01}{\textbf{\textit{#1}}}}
\newcommand{\KeywordTok}[1]{\textcolor[rgb]{0.13,0.29,0.53}{\textbf{#1}}}
\newcommand{\NormalTok}[1]{#1}
\newcommand{\OperatorTok}[1]{\textcolor[rgb]{0.81,0.36,0.00}{\textbf{#1}}}
\newcommand{\OtherTok}[1]{\textcolor[rgb]{0.56,0.35,0.01}{#1}}
\newcommand{\PreprocessorTok}[1]{\textcolor[rgb]{0.56,0.35,0.01}{\textit{#1}}}
\newcommand{\RegionMarkerTok}[1]{#1}
\newcommand{\SpecialCharTok}[1]{\textcolor[rgb]{0.81,0.36,0.00}{\textbf{#1}}}
\newcommand{\SpecialStringTok}[1]{\textcolor[rgb]{0.31,0.60,0.02}{#1}}
\newcommand{\StringTok}[1]{\textcolor[rgb]{0.31,0.60,0.02}{#1}}
\newcommand{\VariableTok}[1]{\textcolor[rgb]{0.00,0.00,0.00}{#1}}
\newcommand{\VerbatimStringTok}[1]{\textcolor[rgb]{0.31,0.60,0.02}{#1}}
\newcommand{\WarningTok}[1]{\textcolor[rgb]{0.56,0.35,0.01}{\textbf{\textit{#1}}}}
\usepackage{graphicx}
\makeatletter
\newsavebox\pandoc@box
\newcommand*\pandocbounded[1]{% scales image to fit in text height/width
  \sbox\pandoc@box{#1}%
  \Gscale@div\@tempa{\textheight}{\dimexpr\ht\pandoc@box+\dp\pandoc@box\relax}%
  \Gscale@div\@tempb{\linewidth}{\wd\pandoc@box}%
  \ifdim\@tempb\p@<\@tempa\p@\let\@tempa\@tempb\fi% select the smaller of both
  \ifdim\@tempa\p@<\p@\scalebox{\@tempa}{\usebox\pandoc@box}%
  \else\usebox{\pandoc@box}%
  \fi%
}
% Set default figure placement to htbp
\def\fps@figure{htbp}
\makeatother
\setlength{\emergencystretch}{3em} % prevent overfull lines
\providecommand{\tightlist}{%
  \setlength{\itemsep}{0pt}\setlength{\parskip}{0pt}}
\usepackage{bookmark}
\IfFileExists{xurl.sty}{\usepackage{xurl}}{} % add URL line breaks if available
\urlstyle{same}
\hypersetup{
  pdftitle={Lab3-hw3},
  hidelinks,
  pdfcreator={LaTeX via pandoc}}

\title{Lab3-hw3}
\author{}
\date{\vspace{-2.5em}}

\begin{document}
\maketitle

\subsection{Overview}\label{overview}

In this lab we'll collect posts from Bluesky Social using the
\texttt{\{atrrr\}} package and apply lexicon-based sentiment analysis
tools --- the same tools we covered in lecture (VADER via
\texttt{\{vader\}}, \texttt{\{sentimentr\}}, and \texttt{\{tidytext\}}
lexicons). By the end you'll have a working pipeline from API call to
sentiment scores, and you'll have a concrete sense of where these tools
succeed and where they fall short. We'll walk through some of this in
class with sample data, then you'll be responsible for completing it
with your own data set for the assignment.

\begin{center}\rule{0.5\linewidth}{0.5pt}\end{center}

\subsection{Setup}\label{setup}

\subsubsection{1. Create a Bluesky account and app password (you'll do
this your
assignment)}\label{create-a-bluesky-account-and-app-password-youll-do-this-your-assignment}

\begin{Shaded}
\begin{Highlighting}[]
\CommentTok{\# Add bluesky account information/ credientials in to the \textasciigrave{}.Renviron\textasciigrave{}}
\CommentTok{\# usethis::edit\_r\_environ() }
\end{Highlighting}
\end{Shaded}

\subsubsection{2.Packages}\label{packages}

\begin{Shaded}
\begin{Highlighting}[]
\CommentTok{\#install.packages(c("atrrr", "tidyverse", "tidytext", "sentimentr", "remotes", "textdata"))}

\CommentTok{\# tidyvader isn\textquotesingle{}t on CRAN:}
\CommentTok{\#install.packages("devtools")}
\CommentTok{\#devtools::install\_github("chris31415926535/tidyvader")}
\end{Highlighting}
\end{Shaded}

\begin{Shaded}
\begin{Highlighting}[]
\FunctionTok{library}\NormalTok{(atrrr)}
\FunctionTok{library}\NormalTok{(tidyverse)}
\end{Highlighting}
\end{Shaded}

\begin{verbatim}
## Warning: package 'ggplot2' was built under R version 4.5.2
\end{verbatim}

\begin{verbatim}
## Warning: package 'tibble' was built under R version 4.5.2
\end{verbatim}

\begin{verbatim}
## Warning: package 'tidyr' was built under R version 4.5.2
\end{verbatim}

\begin{verbatim}
## Warning: package 'readr' was built under R version 4.5.2
\end{verbatim}

\begin{verbatim}
## Warning: package 'purrr' was built under R version 4.5.2
\end{verbatim}

\begin{verbatim}
## Warning: package 'dplyr' was built under R version 4.5.2
\end{verbatim}

\begin{verbatim}
## Warning: package 'lubridate' was built under R version 4.5.2
\end{verbatim}

\begin{verbatim}
## -- Attaching core tidyverse packages ------------------------ tidyverse 2.0.0 --
## v dplyr     1.2.1     v readr     2.1.6
## v forcats   1.0.1     v stringr   1.6.0
## v ggplot2   4.0.2     v tibble    3.3.1
## v lubridate 1.9.5     v tidyr     1.3.2
## v purrr     1.2.1     
## -- Conflicts ------------------------------------------ tidyverse_conflicts() --
## x dplyr::filter() masks stats::filter()
## x dplyr::lag()    masks stats::lag()
## i Use the conflicted package (<http://conflicted.r-lib.org/>) to force all conflicts to become errors
\end{verbatim}

\begin{Shaded}
\begin{Highlighting}[]
\FunctionTok{library}\NormalTok{(tidytext)}
\FunctionTok{library}\NormalTok{(sentimentr)}
\FunctionTok{library}\NormalTok{(remotes)}
\FunctionTok{library}\NormalTok{(textdata)}
\end{Highlighting}
\end{Shaded}

\begin{center}\rule{0.5\linewidth}{0.5pt}\end{center}

\subsection{Part 1: Authenticate and Pull
Posts}\label{part-1-authenticate-and-pull-posts}

\begin{Shaded}
\begin{Highlighting}[]
\CommentTok{\# Authenticate}
\CommentTok{\# Catches token so wont\textquotesingle{}t get hit login rate limits }
\CommentTok{\# auth(}
\CommentTok{\#   user = Sys.getenv("BSKY\_USER"),}
\CommentTok{\#   pass = Sys.getenv("BSKY\_PASS")}
\CommentTok{\#   ) }

\CommentTok{\# Input the term and limit }
\NormalTok{raw\_posts }\OtherTok{\textless{}{-}} \FunctionTok{search\_post}\NormalTok{(}\StringTok{"ecosystem service"}\NormalTok{, }\AttributeTok{limit =} \DecValTok{1000}\NormalTok{) }\CommentTok{\#you\textquotesingle{}ll do this later}
\end{Highlighting}
\end{Shaded}

\begin{verbatim}
## i Found an *unknown total* of posts that fit the query
\end{verbatim}

\begin{verbatim}
## \ Got 100 posts, but there is more.. [742ms]| Got 199 posts, but there is more.. [1s]   / Got 298 posts, but there is more.. [1.3s]- Got 396 posts, but there is more.. [1.5s]\ Got 494 posts, but there is more.. [1.8s]| Got 593 posts, but there is more.. [2.1s]/ Got 688 posts, but there is more.. [2.3s]- Got 786 posts, but there is more.. [2.6s]\ Got 879 posts, but there is more.. [2.8s]| Got 974 posts, but there is more.. [3s]  / Got 1067 posts, but there is more.. [3.2s]                                             i Parsing 1067 results.v Got 1067 results. All done!
\end{verbatim}

Investigate results!

\begin{Shaded}
\begin{Highlighting}[]
\FunctionTok{glimpse}\NormalTok{(raw\_posts)}
\end{Highlighting}
\end{Shaded}

\begin{verbatim}
## Rows: 1,067
## Columns: 21
## $ uri           <chr> "at://did:plc:xhrtfphu4gzqrztkpb3a63qf/app.bsky.feed.pos~
## $ cid           <chr> "bafyreibgbupaqkpf2h2dgavbirplszltiok7l7l5phenpd5zorvyon~
## $ author_handle <chr> "babygoldie.bsky.social", "rwatimes.bsky.social", "maryp~
## $ author_name   <chr> "News by babygoldie", "RWATimes.io", "Mary Branscombe", ~
## $ text          <chr> "[4/5] service capabilities of intermediaries. Promoting~
## $ author_data   <list> ["did:plc:xhrtfphu4gzqrztkpb3a63qf", "babygoldie.bsky.s~
## $ post_data     <list> ["app.bsky.feed.post", "2026-04-29T02:41:15.447303+00:0~
## $ embed_data    <list> <NULL>, <NULL>, <NULL>, <NULL>, <NULL>, <NULL>, <NULL>,~
## $ reply_count   <int> 1, 0, 1, 0, 1, 0, 1, 0, 0, 0, 1, 0, 0, 0, 0, 0, 1, 0, 1,~
## $ repost_count  <int> 0, 0, 0, 0, 0, 0, 0, 0, 0, 0, 0, 0, 0, 0, 0, 0, 0, 0, 0,~
## $ like_count    <int> 0, 0, 1, 0, 0, 0, 0, 0, 0, 2, 0, 0, 0, 1, 0, 0, 0, 0, 0,~
## $ quote_count   <int> 0, 0, 0, 0, 0, 0, 0, 0, 0, 0, 0, 0, 0, 0, 0, 0, 0, 0, 0,~
## $ indexed_at    <dttm> 2026-04-29 02:41:15, 2026-04-28 18:42:12, 2026-04-28 18~
## $ in_reply_to   <chr> "at://did:plc:xhrtfphu4gzqrztkpb3a63qf/app.bsky.feed.pos~
## $ in_reply_root <chr> "at://did:plc:xhrtfphu4gzqrztkpb3a63qf/app.bsky.feed.pos~
## $ quotes        <chr> NA, NA, NA, NA, NA, NA, NA, NA, NA, NA, NA, NA, NA, NA, ~
## $ tags          <list> <NULL>, <NULL>, <NULL>, <NULL>, <NULL>, <NULL>, <NULL>,~
## $ mentions      <list> <NULL>, <NULL>, <NULL>, <NULL>, <NULL>, <NULL>, <NULL>,~
## $ links         <list> <NULL>, <NULL>, <NULL>, <NULL>, <NULL>, <NULL>, <NULL>,~
## $ langs         <list> ["en"], <NULL>, ["en"], ["en"], ["en"], ["en"], ["en"],~
## $ labels        <list> [], [], [], [], [], [], [], [], [], [], [], [], [], [],~
\end{verbatim}

\begin{Shaded}
\begin{Highlighting}[]
\FunctionTok{head}\NormalTok{(raw\_posts)}
\end{Highlighting}
\end{Shaded}

\begin{verbatim}
## # A tibble: 6 x 21
##   uri            cid   author_handle author_name text  author_data  post_data   
##   <chr>          <chr> <chr>         <chr>       <chr> <list>       <list>      
## 1 at://did:plc:~ bafy~ babygoldie.b~ News by ba~ [4/5~ <named list> <named list>
## 2 at://did:plc:~ bafy~ rwatimes.bsk~ RWATimes.io ➤ Th~ <named list> <named list>
## 3 at://did:plc:~ bafy~ marypcbuk.bs~ Mary Brans~ so y~ <named list> <named list>
## 4 at://did:plc:~ bafy~ latticeminds~ Lattice Mi~ Fore~ <named list> <named list>
## 5 at://did:plc:~ bafy~ babygoldie.b~ News by ba~ [5/8~ <named list> <named list>
## 6 at://did:plc:~ bafy~ babygoldie.b~ News by ba~ [6/6~ <named list> <named list>
## # i 14 more variables: embed_data <list>, reply_count <int>,
## #   repost_count <int>, like_count <int>, quote_count <int>, indexed_at <dttm>,
## #   in_reply_to <chr>, in_reply_root <chr>, quotes <chr>, tags <list>,
## #   mentions <list>, links <list>, langs <list>, labels <list>
\end{verbatim}

\begin{center}\rule{0.5\linewidth}{0.5pt}\end{center}

\subsection{Part 2: Clean and Prepare
Text}\label{part-2-clean-and-prepare-text}

Pull out necessary columns \& basic cleaning

\begin{Shaded}
\begin{Highlighting}[]
\NormalTok{posts }\OtherTok{\textless{}{-}}\NormalTok{ raw\_posts }\SpecialCharTok{\%\textgreater{}\%}
  \FunctionTok{select}\NormalTok{(uri, indexed\_at, text, like\_count, reply\_count, repost\_count) }\SpecialCharTok{\%\textgreater{}\%}
  \FunctionTok{mutate}\NormalTok{(}
    \AttributeTok{indexed\_at =} \FunctionTok{as.POSIXct}\NormalTok{(indexed\_at), }\CommentTok{\# who can tell me how to }
    \AttributeTok{post\_date  =} \FunctionTok{as.Date}\NormalTok{(indexed\_at),}
    \AttributeTok{text\_clean =} \FunctionTok{str\_remove\_all}\NormalTok{(text, }\StringTok{"https?}\SpecialCharTok{\textbackslash{}\textbackslash{}}\StringTok{S+"}\NormalTok{), }\CommentTok{\# remove URLs}
    \AttributeTok{text\_clean =} \FunctionTok{str\_remove\_all}\NormalTok{(text\_clean, }\StringTok{"@}\SpecialCharTok{\textbackslash{}\textbackslash{}}\StringTok{S+"}\NormalTok{),}\CommentTok{\# remove mentions}
    \AttributeTok{text\_clean =} \FunctionTok{str\_replace\_all}\NormalTok{(text\_clean, }\StringTok{"\textbackslash{}u2019"}\NormalTok{, }\StringTok{"\textquotesingle{}"}\NormalTok{), }\CommentTok{\# normalize all curly \textquotesingle{} to R\textquotesingle{}s \textquotesingle{}}
    \AttributeTok{text\_clean =} \FunctionTok{str\_remove\_all}\NormalTok{(text\_clean, }\StringTok{"\textquotesingle{}s}\SpecialCharTok{\textbackslash{}\textbackslash{}}\StringTok{b"}\NormalTok{), }\CommentTok{\# remove possessives}
    \AttributeTok{text\_clean =} \FunctionTok{str\_squish}\NormalTok{(text\_clean) }\CommentTok{\# cleans all the added white space }
\NormalTok{  ) }\SpecialCharTok{\%\textgreater{}\%}
  \FunctionTok{filter}\NormalTok{(}\FunctionTok{nchar}\NormalTok{(text\_clean) }\SpecialCharTok{\textgreater{}} \DecValTok{10}\NormalTok{)}
\end{Highlighting}
\end{Shaded}

Number of posts:

\begin{Shaded}
\begin{Highlighting}[]
\FunctionTok{nrow}\NormalTok{(posts) }
\end{Highlighting}
\end{Shaded}

\begin{verbatim}
## [1] 1060
\end{verbatim}

\begin{center}\rule{0.5\linewidth}{0.5pt}\end{center}

\subsection{Part 3: Word Frequency}\label{part-3-word-frequency}

Explore and plot more frequent words used

\begin{Shaded}
\begin{Highlighting}[]
\CommentTok{\# Standard stop words}
\FunctionTok{data}\NormalTok{(}\StringTok{"stop\_words"}\NormalTok{)}

\NormalTok{tokens\_clean }\OtherTok{\textless{}{-}}\NormalTok{ posts }\SpecialCharTok{|\textgreater{}}
  \FunctionTok{unnest\_tokens}\NormalTok{(word, text\_clean) }\SpecialCharTok{|\textgreater{}}
  \FunctionTok{anti\_join}\NormalTok{(stop\_words, }\AttributeTok{by =} \StringTok{"word"}\NormalTok{)}

\CommentTok{\# Custom stop words }
\NormalTok{custom\_sw }\OtherTok{\textless{}{-}} \FunctionTok{tibble}\NormalTok{(}\AttributeTok{word =} \FunctionTok{c}\NormalTok{(}
  \CommentTok{\# search terms}
  \StringTok{"ecosystem"}\NormalTok{, }\StringTok{"service"}
  \CommentTok{\# other common but uninformative text }
  \CommentTok{\# I can\textquotesingle{}t think of anything}
\NormalTok{))}

\NormalTok{tokens\_clean2 }\OtherTok{\textless{}{-}}\NormalTok{ tokens\_clean }\SpecialCharTok{\%\textgreater{}\%}
  \FunctionTok{anti\_join}\NormalTok{(custom\_sw, }\AttributeTok{by =} \StringTok{"word"}\NormalTok{)}

\CommentTok{\# Word frequency plot}

\NormalTok{word\_counts }\OtherTok{\textless{}{-}}\NormalTok{ tokens\_clean2 }\SpecialCharTok{\%\textgreater{}\%} 
  \FunctionTok{count}\NormalTok{(word, }\AttributeTok{sort =} \ConstantTok{TRUE}\NormalTok{)}

\NormalTok{word\_counts }\SpecialCharTok{\%\textgreater{}\%}
  \FunctionTok{slice\_max}\NormalTok{(n, }\AttributeTok{n =}\DecValTok{20}\NormalTok{) }\SpecialCharTok{\%\textgreater{}\%} 
  \FunctionTok{ggplot}\NormalTok{(}\FunctionTok{aes}\NormalTok{(}\AttributeTok{x =} \FunctionTok{reorder}\NormalTok{(word, n), }\AttributeTok{y =}\NormalTok{ n)) }\SpecialCharTok{+} 
  \FunctionTok{geom\_col}\NormalTok{(}\AttributeTok{fill =} \StringTok{"\#6c91c9"}\NormalTok{) }\SpecialCharTok{+} 
  \FunctionTok{coord\_flip}\NormalTok{() }\SpecialCharTok{+} 
  \FunctionTok{labs}\NormalTok{(}\AttributeTok{title =} \StringTok{"top 20 words in Blue Sky \textquotesingle{}Ecosystem Service\textquotesingle{} posts"}\NormalTok{) }\SpecialCharTok{+} 
  \FunctionTok{theme\_minimal}\NormalTok{()}
\end{Highlighting}
\end{Shaded}

\pandocbounded{\includegraphics[keepaspectratio]{Lab3-hw3_files/Lab3-hw3_files/figure-latex/word-freq-1.pdf}}

\subsection{Part 4: Sentiment Analysis}\label{part-4-sentiment-analysis}

\begin{quote}
\textbf{Note:} If you need to re-run any chunk in this section, run
\texttt{rm(posts)} in the console first and re-run from Part 2. This
avoids duplicate column errors.
\end{quote}

\subsubsection{4a. Bing lexicon
(word-level)}\label{a.-bing-lexicon-word-level}

The Bing lexicon is a polarity / binary lexicon grouping where yes/no =
positive/negative. It is a list of words and their positive or negative
sentiments.

\begin{Shaded}
\begin{Highlighting}[]
\CommentTok{\# Load lexicon }
\NormalTok{bing\_lexicon }\OtherTok{\textless{}{-}} \FunctionTok{get\_sentiments}\NormalTok{(}\StringTok{"bing"}\NormalTok{)}
\CommentTok{\#get\_sentiments("bing")}
\end{Highlighting}
\end{Shaded}

Join the bing lexicon list to the tokenized the posts and calcultae the
net score (positive word count minus negative word count) per post.

\begin{Shaded}
\begin{Highlighting}[]
\CommentTok{\# Score posts with Bing}
\NormalTok{bing\_scores }\OtherTok{\textless{}{-}}\NormalTok{ posts }\SpecialCharTok{\%\textgreater{}\%} 
  \CommentTok{\# Tokenize posts }
  \FunctionTok{unnest\_tokens}\NormalTok{(word, text\_clean) }\SpecialCharTok{\%\textgreater{}\%} 
  \CommentTok{\# Join the words to the sentiment lexicon }
  \FunctionTok{inner\_join}\NormalTok{(bing\_lexicon, }\AttributeTok{by =} \StringTok{"word"}\NormalTok{) }\SpecialCharTok{\%\textgreater{}\%} 
  \FunctionTok{count}\NormalTok{(uri, sentiment) }\SpecialCharTok{\%\textgreater{}\%} 
  \FunctionTok{pivot\_wider}\NormalTok{(}\AttributeTok{names\_from =}\NormalTok{ sentiment, }
              \AttributeTok{values\_from =}\NormalTok{ n, }
              \AttributeTok{values\_fill =} \DecValTok{0}\NormalTok{) }\SpecialCharTok{\%\textgreater{}\%} 
  \CommentTok{\# Calculate total sentiment score: positives {-} negative }
  \FunctionTok{mutate}\NormalTok{(}\AttributeTok{score\_b =}\NormalTok{ positive }\SpecialCharTok{{-}}\NormalTok{ negative) }
\end{Highlighting}
\end{Shaded}

Join the scores back onto \texttt{posts}

\begin{Shaded}
\begin{Highlighting}[]
\CommentTok{\# Join scores to posts  }
\NormalTok{posts }\OtherTok{\textless{}{-}}\NormalTok{ posts }\SpecialCharTok{\%\textgreater{}\%}
  \FunctionTok{left\_join}\NormalTok{(bing\_scores, }\StringTok{"uri"}\NormalTok{) }\SpecialCharTok{\%\textgreater{}\%} 
  \FunctionTok{mutate}\NormalTok{(}\AttributeTok{bing\_score =} \FunctionTok{replace\_na}\NormalTok{(score\_b, }\DecValTok{0}\NormalTok{)) }\CommentTok{\# NA scores = 0 }

\CommentTok{\# Check if it joined the score column }
\CommentTok{\# colnames(posts)}
\end{Highlighting}
\end{Shaded}

\subsubsection{4b. AFINN lexicon (numeric
scores)}\label{b.-afinn-lexicon-numeric-scores}

Score the posts with the afinn sentiment package where each word carries
an integer value (−5 to +5).

\begin{Shaded}
\begin{Highlighting}[]
\CommentTok{\# You will be prompted to download the AFINN lexicon the first time —}
\CommentTok{\# run this line in the console and select Yes, then it will be cached locally}
\NormalTok{afinn }\OtherTok{\textless{}{-}} \FunctionTok{get\_sentiments}\NormalTok{(}\StringTok{"afinn"}\NormalTok{)}

\CommentTok{\# Score text}
\NormalTok{afinn\_scores }\OtherTok{\textless{}{-}}\NormalTok{ posts }\SpecialCharTok{\%\textgreater{}\%}
  \CommentTok{\# Tokenize posts }
  \FunctionTok{unnest\_tokens}\NormalTok{(word, text\_clean) }\SpecialCharTok{\%\textgreater{}\%} 
  \CommentTok{\# Inner Join the words to the sentiment lexicon  }
  \FunctionTok{inner\_join}\NormalTok{(afinn, }\AttributeTok{by =} \StringTok{"word"}\NormalTok{) }\SpecialCharTok{\%\textgreater{}\%} 
  \CommentTok{\# Score: sum of afinn scores for each URI }
  \FunctionTok{group\_by}\NormalTok{(uri) }\SpecialCharTok{\%\textgreater{}\%} 
  \FunctionTok{summarise}\NormalTok{(}\AttributeTok{score\_af =} \FunctionTok{sum}\NormalTok{(value), }\AttributeTok{.groups =} \StringTok{"drop"}\NormalTok{)}


\CommentTok{\# Join scores to posts  }
\NormalTok{posts }\OtherTok{\textless{}{-}}\NormalTok{ posts }\SpecialCharTok{\%\textgreater{}\%} 
  \FunctionTok{left\_join}\NormalTok{(afinn\_scores, }\StringTok{"uri"}\NormalTok{) }\SpecialCharTok{\%\textgreater{}\%} 
  \FunctionTok{mutate}\NormalTok{(}\AttributeTok{afinn\_score =} \FunctionTok{replace\_na}\NormalTok{(score\_af, }\DecValTok{0}\NormalTok{)) }\CommentTok{\# NA scores = 0 }

\CommentTok{\# str(posts) \# Checking dataset }
\end{Highlighting}
\end{Shaded}

Clean the dataset before moving on

\begin{Shaded}
\begin{Highlighting}[]
\NormalTok{posts }\OtherTok{\textless{}{-}}\NormalTok{ posts }\SpecialCharTok{\%\textgreater{}\%} 
  \FunctionTok{select}\NormalTok{(}\SpecialCharTok{{-}}\FunctionTok{c}\NormalTok{(negative, positive, score\_b, score\_af))}
\end{Highlighting}
\end{Shaded}

\subsubsection{4c. NRC lexicon (emotion
categories)}\label{c.-nrc-lexicon-emotion-categories}

The NRC lexicon tags words with 8 emotions (anger, fear, anticipation,
trust, surprise, sadness, joy, disgust) plus positive/negative.

Now, count emotion words per post and visualize which emotions dominate
the corpus.

\begin{Shaded}
\begin{Highlighting}[]
\NormalTok{nrc }\OtherTok{\textless{}{-}} \FunctionTok{get\_sentiments}\NormalTok{(}\StringTok{"nrc"}\NormalTok{)}

\CommentTok{\# Score text {-} Count emotion words per post}
\NormalTok{nrc\_scores }\OtherTok{\textless{}{-}}\NormalTok{ posts }\SpecialCharTok{\%\textgreater{}\%} 
  \CommentTok{\# Tokenized posts }
  \FunctionTok{unnest\_tokens}\NormalTok{(word, text\_clean) }\SpecialCharTok{\%\textgreater{}\%} 
  \CommentTok{\# Join to NRC }
  \FunctionTok{inner\_join}\NormalTok{(nrc, }\AttributeTok{by =} \StringTok{"word"}\NormalTok{, }\AttributeTok{relationship =} \StringTok{"many{-}to{-}many"}\NormalTok{) }\SpecialCharTok{\%\textgreater{}\%} 
  \CommentTok{\# Count emotion words per post }
  \FunctionTok{group\_by}\NormalTok{(uri, sentiment) }\SpecialCharTok{\%\textgreater{}\%} 
  \FunctionTok{summarise}\NormalTok{(}\AttributeTok{nrc\_count =} \FunctionTok{n}\NormalTok{(), }\AttributeTok{.groups =} \StringTok{"drop"}\NormalTok{) }


\CommentTok{\# Visualize which emotions dominate the corpus}
\NormalTok{nrc\_scores }\SpecialCharTok{\%\textgreater{}\%}
  \FunctionTok{group\_by}\NormalTok{(sentiment) }\SpecialCharTok{\%\textgreater{}\%}
  \FunctionTok{summarise}\NormalTok{(}\AttributeTok{nrc\_count =} \FunctionTok{sum}\NormalTok{(nrc\_count), }\AttributeTok{.groups =} \StringTok{\textquotesingle{}drop\textquotesingle{}}\NormalTok{) }\SpecialCharTok{\%\textgreater{}\%}
  \FunctionTok{ggplot}\NormalTok{() }\SpecialCharTok{+}
  \FunctionTok{geom\_col}\NormalTok{(}\FunctionTok{aes}\NormalTok{(}\AttributeTok{y =} \FunctionTok{reorder}\NormalTok{(sentiment, nrc\_count), }\AttributeTok{x =}\NormalTok{ nrc\_count), }
           \AttributeTok{fill =} \StringTok{"\#e14225"}\NormalTok{) }\SpecialCharTok{+}
  \FunctionTok{theme\_minimal}\NormalTok{() }\SpecialCharTok{+}
  \FunctionTok{labs}\NormalTok{(}\AttributeTok{title =} \StringTok{"Emotion Words in Search"}\NormalTok{, }\AttributeTok{x =} \StringTok{"Count"}\NormalTok{, }\AttributeTok{y =} \StringTok{"Sentiment Words"}\NormalTok{)}
\end{Highlighting}
\end{Shaded}

\pandocbounded{\includegraphics[keepaspectratio]{Lab3-hw3_files/Lab3-hw3_files/figure-latex/nrc-1.pdf}}

\begin{Shaded}
\begin{Highlighting}[]
\CommentTok{\#Join to post}
\NormalTok{posts }\OtherTok{\textless{}{-}}\NormalTok{ posts }\SpecialCharTok{\%\textgreater{}\%}
  \FunctionTok{left\_join}\NormalTok{(nrc\_scores, }\AttributeTok{by =} \StringTok{"uri"}\NormalTok{) }\SpecialCharTok{\%\textgreater{}\%} 
  \FunctionTok{rename}\NormalTok{(}\AttributeTok{nrc\_sentiment\_words =}\NormalTok{ sentiment)}
\end{Highlighting}
\end{Shaded}

Which emotions dominate the corpus? Plot the count of total scored words
for each of the 8 emotions.

\begin{Shaded}
\begin{Highlighting}[]
\NormalTok{nrc\_scores }\SpecialCharTok{\%\textgreater{}\%}
  \FunctionTok{group\_by}\NormalTok{(sentiment) }\SpecialCharTok{\%\textgreater{}\%}
  \FunctionTok{summarise}\NormalTok{(}\AttributeTok{nrc\_count =} \FunctionTok{sum}\NormalTok{(nrc\_count)) }\SpecialCharTok{\%\textgreater{}\%} 
  \FunctionTok{filter}\NormalTok{(}\SpecialCharTok{!}\NormalTok{sentiment }\SpecialCharTok{\%in\%} \FunctionTok{c}\NormalTok{(}\StringTok{"positive"}\NormalTok{, }\StringTok{"negative"}\NormalTok{)) }\SpecialCharTok{\%\textgreater{}\%} \CommentTok{\# Remove positive and negative }
  \FunctionTok{ggplot}\NormalTok{() }\SpecialCharTok{+}
  \FunctionTok{geom\_col}\NormalTok{(}\FunctionTok{aes}\NormalTok{(}\AttributeTok{y =} \FunctionTok{reorder}\NormalTok{(sentiment, nrc\_count), }\AttributeTok{x =}\NormalTok{ nrc\_count), }\AttributeTok{fill =} \StringTok{"goldenrod"}\NormalTok{) }\SpecialCharTok{+}
  \FunctionTok{theme\_minimal}\NormalTok{() }\SpecialCharTok{+}
  \FunctionTok{labs}\NormalTok{(}\AttributeTok{title =} \StringTok{"Count of Emotion Words in Search"}\NormalTok{, }\AttributeTok{x =} \StringTok{"Count"}\NormalTok{, }\AttributeTok{y =} \StringTok{"Sentiment Words"}\NormalTok{)}
\end{Highlighting}
\end{Shaded}

\pandocbounded{\includegraphics[keepaspectratio]{Lab3-hw3_files/Lab3-hw3_files/figure-latex/nrc-plot-1.pdf}}

What words are driving each emotion? Create plots of the top 10 words
for each emotion (hint: ggplot and
\texttt{facet\_wrap(\textasciitilde{}sentiment)}).

\begin{Shaded}
\begin{Highlighting}[]
\CommentTok{\# Find frequency of words for each sentiment }
\NormalTok{sentiment\_word\_count }\OtherTok{\textless{}{-}}\NormalTok{ posts }\SpecialCharTok{\%\textgreater{}\%} 
  \FunctionTok{unnest\_tokens}\NormalTok{(word, text\_clean) }\SpecialCharTok{\%\textgreater{}\%} \CommentTok{\# Tokenized posts}
  \FunctionTok{inner\_join}\NormalTok{(nrc, }\AttributeTok{by =} \StringTok{"word"}\NormalTok{, }\AttributeTok{relationship =} \StringTok{"many{-}to{-}many"}\NormalTok{) }\SpecialCharTok{\%\textgreater{}\%} \CommentTok{\# Join to NRC }
  \FunctionTok{count}\NormalTok{(sentiment, word, }\AttributeTok{sort =} \ConstantTok{TRUE}\NormalTok{)  }\CommentTok{\# Count of words per sentiment }

\CommentTok{\# Get the top 10 most frequent words for each sentiment group }
\NormalTok{top\_10 }\OtherTok{\textless{}{-}}\NormalTok{ sentiment\_word\_count }\SpecialCharTok{\%\textgreater{}\%} 
  \FunctionTok{group\_by}\NormalTok{(sentiment) }\SpecialCharTok{\%\textgreater{}\%} 
  \FunctionTok{slice\_max}\NormalTok{(}\AttributeTok{order\_by =}\NormalTok{ n, }\AttributeTok{n =} \DecValTok{10}\NormalTok{, }\AttributeTok{with\_ties =} \ConstantTok{FALSE}\NormalTok{) }\SpecialCharTok{\%\textgreater{}\%}
  \FunctionTok{ungroup}\NormalTok{()}

\CommentTok{\# Plot the top 10 words for each sentiment group }
\FunctionTok{ggplot}\NormalTok{(top\_10, }\FunctionTok{aes}\NormalTok{(}\AttributeTok{x =} \FunctionTok{reorder\_within}\NormalTok{(word, n, sentiment), }\AttributeTok{y =}\NormalTok{ n, }\AttributeTok{fill =}\NormalTok{ sentiment)) }\SpecialCharTok{+}
  \FunctionTok{geom\_col}\NormalTok{(}\AttributeTok{show.legend =} \ConstantTok{FALSE}\NormalTok{) }\SpecialCharTok{+}
  \FunctionTok{scale\_x\_reordered}\NormalTok{() }\SpecialCharTok{+}  \CommentTok{\# Keeps ordering within facets}
  \FunctionTok{facet\_wrap}\NormalTok{(}\SpecialCharTok{\textasciitilde{}}\NormalTok{ sentiment, }\AttributeTok{scales =} \StringTok{"free\_y"}\NormalTok{) }\SpecialCharTok{+}
  \FunctionTok{labs}\NormalTok{(}
    \AttributeTok{title =} \StringTok{"Top 10 Words per Sentiment Group"}\NormalTok{,}
    \AttributeTok{x =} \StringTok{"Word"}\NormalTok{,}
    \AttributeTok{y =} \StringTok{"Frequency"}
\NormalTok{  ) }\SpecialCharTok{+}
  \FunctionTok{coord\_flip}\NormalTok{() }\SpecialCharTok{+} 
  \FunctionTok{theme\_minimal}\NormalTok{()}
\end{Highlighting}
\end{Shaded}

\pandocbounded{\includegraphics[keepaspectratio]{Lab3-hw3_files/Lab3-hw3_files/figure-latex/nrc-top-words-1.pdf}}

\subsubsection{4d. sentimentr
(sentence-aware)}\label{d.-sentimentr-sentence-aware}

Unlike the word-level lexicons above, \texttt{\{sentimentr\}} scores at
the sentence level and handles valence shifters like negation --- ``not
alarming'' won't score the same as ``alarming.'' Pass the full
\texttt{text\_clean} vector to \texttt{sentiment\_by()}. Rename the
resulting \texttt{ave\_sentiment} to \texttt{sr\_scores} and join back
onto \texttt{posts} by row index. I'll give you this code as we haven't
fully covered this yet.

\begin{Shaded}
\begin{Highlighting}[]
\CommentTok{\# Using sentimentr }
\NormalTok{sr\_scores }\OtherTok{\textless{}{-}} \FunctionTok{sentiment\_by}\NormalTok{(posts}\SpecialCharTok{$}\NormalTok{text\_clean) }\SpecialCharTok{\%\textgreater{}\%}
  \FunctionTok{select}\NormalTok{(element\_id, ave\_sentiment) }\SpecialCharTok{\%\textgreater{}\%}
  \FunctionTok{rename}\NormalTok{(}\AttributeTok{sr\_score =}\NormalTok{ ave\_sentiment) }\SpecialCharTok{\%\textgreater{}\%}
  \FunctionTok{mutate}\NormalTok{(}\AttributeTok{element\_id =} \FunctionTok{as.integer}\NormalTok{(element\_id))}
\end{Highlighting}
\end{Shaded}

\begin{verbatim}
## Warning: Each time `sentiment_by` is run it has to do sentence boundary disambiguation when a
## raw `character` vector is passed to `text.var`. This may be costly of time and
## memory.  It is highly recommended that the user first runs the raw `character`
## vector through the `get_sentences` function.
\end{verbatim}

\begin{Shaded}
\begin{Highlighting}[]
\CommentTok{\# Join sr score back into posts by row index }
\NormalTok{posts }\OtherTok{\textless{}{-}}\NormalTok{ posts }\SpecialCharTok{\%\textgreater{}\%}
  \FunctionTok{mutate}\NormalTok{(}\AttributeTok{element\_id =} \FunctionTok{row\_number}\NormalTok{()) }\SpecialCharTok{\%\textgreater{}\%}
  \FunctionTok{left\_join}\NormalTok{(sr\_scores, }\AttributeTok{by =} \StringTok{"element\_id"}\NormalTok{)}
\end{Highlighting}
\end{Shaded}

\subsubsection{4e. VADER}\label{e.-vader}

VADER scores the full post text (no tokenization needed) and returns
four values: \texttt{pos}, \texttt{neu}, \texttt{neg}, and
\texttt{compound}. Use \texttt{tidyvader::vader()} on the
\texttt{text\_clean} column --- it mutates the result directly onto the
dataframe. Keep only the \texttt{compound} score and drop the component
columns.

Let's run VADER on a single post manually before applying it to the
whole dataset:

\begin{Shaded}
\begin{Highlighting}[]
\CommentTok{\# look at a few posts}
\NormalTok{posts }\SpecialCharTok{\%\textgreater{}\%} \FunctionTok{select}\NormalTok{(text\_clean) }\SpecialCharTok{\%\textgreater{}\%} \FunctionTok{slice}\NormalTok{(}\DecValTok{1}\SpecialCharTok{:}\DecValTok{10}\NormalTok{)}
\end{Highlighting}
\end{Shaded}

\begin{verbatim}
## # A tibble: 10 x 1
##    text_clean                                                                   
##    <chr>                                                                        
##  1 [4/5] service capabilities of intermediaries. Promoting the high-quality dev~
##  2 [4/5] service capabilities of intermediaries. Promoting the high-quality dev~
##  3 [4/5] service capabilities of intermediaries. Promoting the high-quality dev~
##  4 [4/5] service capabilities of intermediaries. Promoting the high-quality dev~
##  5 [4/5] service capabilities of intermediaries. Promoting the high-quality dev~
##  6 ➤ The card integrates with Tether digital wallet, expanding its ecosystem in~
##  7 ➤ The card integrates with Tether digital wallet, expanding its ecosystem in~
##  8 ➤ The card integrates with Tether digital wallet, expanding its ecosystem in~
##  9 so you've got the direct impact on what people buy through and from Amazon a~
## 10 so you've got the direct impact on what people buy through and from Amazon a~
\end{verbatim}

\begin{Shaded}
\begin{Highlighting}[]
\CommentTok{\# run VADER on one post — swap in whichever row looks interesting to you}
\NormalTok{posts}\SpecialCharTok{$}\NormalTok{text\_clean[}\DecValTok{3}\NormalTok{]}
\end{Highlighting}
\end{Shaded}

\begin{verbatim}
## [1] "[4/5] service capabilities of intermediaries. Promoting the high-quality development of listed companies is a systematic project. All regions and departments must take the initiative to optimize the business environment, jointly create a good ecosystem that supports the high-quality"
\end{verbatim}

\begin{Shaded}
\begin{Highlighting}[]
\NormalTok{tidyvader}\SpecialCharTok{::}\FunctionTok{vader\_chr}\NormalTok{(posts}\SpecialCharTok{$}\NormalTok{text\_clean[}\DecValTok{3}\NormalTok{])}
\end{Highlighting}
\end{Shaded}

\begin{verbatim}
##   compound    pos    neu neg
## 1   0.9042 0.2857 0.7143   0
\end{verbatim}

What did VADER catch? What did it miss? How does the compound score
relate to the word-level scores?

VADAR determined the whole text was neutral (neu = 1), which makes sense
as the sentence is a factual, descriptive statement with no emotional
language.

\subsection{\texorpdfstring{Now score all the posts with VADER. Rename
it's default scoring column to
\texttt{vader\_score}.}{Now score all the posts with VADER. Rename it's default scoring column to vader\_score.}}\label{now-score-all-the-posts-with-vader.-rename-its-default-scoring-column-to-vader_score.}

\begin{Shaded}
\begin{Highlighting}[]
\NormalTok{posts }\OtherTok{\textless{}{-}}\NormalTok{ posts }\SpecialCharTok{\%\textgreater{}\%} 
  \FunctionTok{mutate}\NormalTok{(}\AttributeTok{vader\_score =}\NormalTok{ purrr}\SpecialCharTok{::}\FunctionTok{map\_dbl}\NormalTok{(text\_clean, }\SpecialCharTok{\textasciitilde{}}\NormalTok{ tidyvader}\SpecialCharTok{::}\FunctionTok{vader\_chr}\NormalTok{(.x)}\SpecialCharTok{$}\NormalTok{compound))}
\end{Highlighting}
\end{Shaded}

\begin{center}\rule{0.5\linewidth}{0.5pt}\end{center}

\subsection{Part 5: Compare the
Methods}\label{part-5-compare-the-methods}

\subsubsection{Correlations}\label{correlations}

Let's see to what extent the 4 numerical methods (bing, afinn, sr, and
vader) are in agreement. Select the 4 scores and use \texttt{cor()}
using ``complete.obs''.

\begin{Shaded}
\begin{Highlighting}[]
\NormalTok{posts }\SpecialCharTok{\%\textgreater{}\%} 
  \FunctionTok{select}\NormalTok{(bing\_score, afinn\_score, sr\_score, vader\_score) }\SpecialCharTok{\%\textgreater{}\%} \CommentTok{\# Select 4 scores }
  \FunctionTok{cor}\NormalTok{(}\AttributeTok{use =} \StringTok{"complete.obs"}\NormalTok{) }\CommentTok{\# use cor()}
\end{Highlighting}
\end{Shaded}

\begin{verbatim}
##             bing_score afinn_score  sr_score vader_score
## bing_score   1.0000000   0.6931579 0.6006261   0.6634219
## afinn_score  0.6931579   1.0000000 0.5079764   0.7735428
## sr_score     0.6006261   0.5079764 1.0000000   0.5563381
## vader_score  0.6634219   0.7735428 0.5563381   1.0000000
\end{verbatim}

\subsubsection{Score distributions}\label{score-distributions}

Use \texttt{pivot\_longer()} to reshape the four score columns into a
single \texttt{score} column with a \texttt{method} identifier, then
plot with \texttt{geom\_histogram()} and \texttt{facet\_wrap()}.

\begin{Shaded}
\begin{Highlighting}[]
\NormalTok{posts }\SpecialCharTok{\%\textgreater{}\%} 
  \FunctionTok{pivot\_longer}\NormalTok{(}\AttributeTok{cols =} \FunctionTok{c}\NormalTok{(bing\_score, afinn\_score, sr\_score, vader\_score), }
               \AttributeTok{names\_to =} \StringTok{"method"}\NormalTok{, }
               \AttributeTok{values\_to =} \StringTok{"scores"}\NormalTok{) }\SpecialCharTok{\%\textgreater{}\%} 
  \FunctionTok{ggplot}\NormalTok{() }\SpecialCharTok{+} 
  \FunctionTok{geom\_histogram}\NormalTok{(}\FunctionTok{aes}\NormalTok{(scores, }\AttributeTok{fill =}\NormalTok{ method), }\AttributeTok{show.legend =} \ConstantTok{FALSE}\NormalTok{) }\SpecialCharTok{+} 
  \FunctionTok{facet\_wrap}\NormalTok{(}\SpecialCharTok{\textasciitilde{}}\NormalTok{method) }\SpecialCharTok{+} 
  \FunctionTok{theme\_bw}\NormalTok{()}
\end{Highlighting}
\end{Shaded}

\begin{verbatim}
## `stat_bin()` using `bins = 30`. Pick better value `binwidth`.
\end{verbatim}

\pandocbounded{\includegraphics[keepaspectratio]{Lab3-hw3_files/Lab3-hw3_files/figure-latex/distributions-1.pdf}}

\textbf{Question 1:} Do the four methods broadly agree or disagree?
Where do you see the biggest differences in the distributions? Which
method produces the most zeros, and why?

All four methods show a neutral corpus.

The VADAR and sentimentr results are similar. Their scores are very
neutral, having high counts at/around 0. These are both sentence-level
analyzers that account for context. These methods are more precise but
more conservative as they miss sentiment that is not explicitly
emotional. The accumulated text does not contain strong language, so
VADAR and sentimentr scores are concentrated at 0.

Unlike VADAR and sentimentr, the bing and afinn scores both have a large
score range and broader distributions. Bing and afinn are word-count
based so they accumulate scores across a document, resulting in more
variations across posts.

\begin{center}\rule{0.5\linewidth}{0.5pt}\end{center}

\subsection{Part 6: Spot-Check for
Misclassification}\label{part-6-spot-check-for-misclassification}

Find posts where the methods disagree most --- these are the best
candidates for understanding failure modes. Use \texttt{mutate()} to
create binary flags from the continuous scores, then \texttt{filter()}
for posts where the two methods point in opposite directions.

Finding where bing and VADAD scores point opposite directions.

\begin{Shaded}
\begin{Highlighting}[]
\NormalTok{misclassifications }\OtherTok{\textless{}{-}}\NormalTok{ posts }\SpecialCharTok{\%\textgreater{}\%} 
  \FunctionTok{select}\NormalTok{(}\FunctionTok{c}\NormalTok{(text\_clean, bing\_score, vader\_score)) }\SpecialCharTok{\%\textgreater{}\%} 
  \CommentTok{\# New column with score signs }
  \FunctionTok{mutate}\NormalTok{(}\AttributeTok{bing\_flag =} \FunctionTok{sign}\NormalTok{(bing\_score),}
         \CommentTok{\#afinn\_flag = sign(afinn\_score),}
         \CommentTok{\#sr\_flag = sign(sr\_score),}
         \AttributeTok{vadar\_flag =} \FunctionTok{sign}\NormalTok{(vader\_score)}
\NormalTok{         ) }\SpecialCharTok{\%\textgreater{}\%} 
  \CommentTok{\# Filter for posts where two methods point in opposite directions }
  \CommentTok{\# Picking bing and VADAR }
  \FunctionTok{filter}\NormalTok{((bing\_flag }\SpecialCharTok{==} \SpecialCharTok{{-}}\DecValTok{1} \SpecialCharTok{\&}\NormalTok{ vadar\_flag }\SpecialCharTok{==} \DecValTok{1}\NormalTok{) }\SpecialCharTok{|} 
\NormalTok{         (bing\_flag }\SpecialCharTok{==} \DecValTok{1} \SpecialCharTok{\&}\NormalTok{ vadar\_flag }\SpecialCharTok{==} \SpecialCharTok{{-}}\DecValTok{1}\NormalTok{))}

\FunctionTok{head}\NormalTok{(misclassifications)}
\end{Highlighting}
\end{Shaded}

\begin{verbatim}
## # A tibble: 6 x 5
##   text_clean                         bing_score vader_score bing_flag vadar_flag
##   <chr>                                   <int>       <dbl>     <dbl>      <dbl>
## 1 I'm not sure it makes up losses l~          1      -0.287         1         -1
## 2 I'm not sure it makes up losses l~          1      -0.287         1         -1
## 3 I'm not sure it makes up losses l~          1      -0.287         1         -1
## 4 Basically use Gamepass as a promo~         -1       0.226        -1          1
## 5 Basically use Gamepass as a promo~         -1       0.226        -1          1
## 6 Basically use Gamepass as a promo~         -1       0.226        -1          1
\end{verbatim}

To rank posts by total disagreement, compute the absolute difference
between the two scores with \texttt{abs()}, then
\texttt{arrange(desc(...))} to surface the highest-conflict posts.

\begin{Shaded}
\begin{Highlighting}[]
\NormalTok{high\_conflict\_text }\OtherTok{\textless{}{-}}\NormalTok{ posts }\SpecialCharTok{\%\textgreater{}\%} 
  \CommentTok{\# Text per row and keep only neccessary columns  }
  \FunctionTok{select}\NormalTok{(}\FunctionTok{c}\NormalTok{(text\_clean, bing\_score, vader\_score)) }\SpecialCharTok{\%\textgreater{}\%} 
  \FunctionTok{group\_by}\NormalTok{(text\_clean) }\SpecialCharTok{\%\textgreater{}\%} 
  \FunctionTok{summarise}\NormalTok{(}\AttributeTok{bing\_score =} \FunctionTok{mean}\NormalTok{(bing\_score), }
            \AttributeTok{vader\_score =} \FunctionTok{mean}\NormalTok{(vader\_score)) }\SpecialCharTok{\%\textgreater{}\%} 
  \CommentTok{\# New column with score signs }
  \FunctionTok{mutate}\NormalTok{(}\AttributeTok{bing\_flag =} \FunctionTok{sign}\NormalTok{(bing\_score),}
         \CommentTok{\#afinn\_flag = sign(afinn\_score),}
         \CommentTok{\#sr\_flag = sign(sr\_score),}
         \AttributeTok{vadar\_flag =} \FunctionTok{sign}\NormalTok{(vader\_score)}
\NormalTok{         ) }\SpecialCharTok{\%\textgreater{}\%} 
  \CommentTok{\# Filter for posts where two methods point in opposite directions }
  \CommentTok{\# Picking bing and VADAR }
  \FunctionTok{filter}\NormalTok{((bing\_flag }\SpecialCharTok{==} \SpecialCharTok{{-}}\DecValTok{1} \SpecialCharTok{\&}\NormalTok{ vadar\_flag }\SpecialCharTok{==} \DecValTok{1}\NormalTok{) }\SpecialCharTok{|} 
\NormalTok{         (bing\_flag }\SpecialCharTok{==} \DecValTok{1} \SpecialCharTok{\&}\NormalTok{ vadar\_flag }\SpecialCharTok{==} \SpecialCharTok{{-}}\DecValTok{1}\NormalTok{)) }\SpecialCharTok{\%\textgreater{}\%} 
  \CommentTok{\# Compute the absolute difference between the two scores }
  \FunctionTok{mutate}\NormalTok{(}\AttributeTok{abs\_difference =} \FunctionTok{abs}\NormalTok{(bing\_score) }\SpecialCharTok{{-}} \FunctionTok{abs}\NormalTok{(vader\_score)) }\SpecialCharTok{\%\textgreater{}\%} 
  \CommentTok{\# Arrange for highest{-}conflict first }
  \FunctionTok{arrange}\NormalTok{(}\FunctionTok{desc}\NormalTok{(abs\_difference))}

\FunctionTok{head}\NormalTok{(high\_conflict\_text)}
\end{Highlighting}
\end{Shaded}

\begin{verbatim}
## # A tibble: 6 x 6
##   text_clean          bing_score vader_score bing_flag vadar_flag abs_difference
##   <chr>                    <dbl>       <dbl>     <dbl>      <dbl>          <dbl>
## 1 Its an ecosystem. ~          6     -0.151          1         -1           5.85
## 2 14) Kosher wedding~          3     -0.0687         1         -1           2.93
## 3 this suspect n/c i~         -3      0.209         -1          1           2.79
## 4 Urgent security up~         -3      0.310         -1          1           2.69
## 5 tanstack.com/ai/la~          3     -0.422          1         -1           2.58
## 6 Honestly shocked a~         -3      0.643         -1          1           2.36
\end{verbatim}

\textbf{Question 2:} Read through 5 of the high-conflict posts. For each
one, write a sentence explaining which method you think is closer to
correct and why.

\textbf{Row 1}: Bing is closer. The post is frustrated/critical but Bing
catches the positive framing around the industry. VADER's negative score
misses that the negativity is directed at policymakers, not the subject.

\textbf{Row 2}: Bing is closer. Mixed sentiment, net positive is
reasonable. VADER going negative likely triggered on ``shit'' without
capturing the overall constructive tone.

\textbf{Row 3}: VADER is closer. The post is clearly critical/negative.
Bing scoring +3 likely matched words like ``proud,'' ``massive,''
``peacemaker'' without understanding they're used sarcastically or
negatively.

\textbf{Row 4}: Neither is great, but VADER is wrong. Scoring +0.31 on
an urgent security warning is clearly off. Bing's -3 at least gets the
direction right, likely from ``vulnerability,'' ``flaw,'' ``denial.''

\textbf{Row 5}: Bing is closer. The post is promotional/positive. VADER
scoring -0.42 likely penalized ``no middleman, no fees, no lock-in''
treating the negations as negative sentiment rather than positive
selling points.

\subsection{\texorpdfstring{\textbf{Question 3:} Based on everything
you've run, which method would you choose if you were delivering a
report onpublic sentiment around your topic, and why? What additional
steps would you take to validate your
choice?}{Question 3: Based on everything you've run, which method would you choose if you were delivering a report onpublic sentiment around your topic, and why? What additional steps would you take to validate your choice?}}\label{question-3-based-on-everything-youve-run-which-method-would-you-choose-if-you-were-delivering-a-report-onpublic-sentiment-around-your-topic-and-why-what-additional-steps-would-you-take-to-validate-your-choice}

\subsection{Bonus: Engagement vs.~Sentiment (completion
optional)}\label{bonus-engagement-vs.-sentiment-completion-optional}

Do more emotionally charged posts get more engagement? Use
\texttt{abs()} on the VADER compound score to capture emotional
intensity regardless of direction, apply \texttt{log1p()} to
\texttt{like\_count} to handle the skewed distribution, and add
\texttt{geom\_smooth(method\ =\ "lm")} to visualize the trend.

\begin{Shaded}
\begin{Highlighting}[]
\NormalTok{posts }\SpecialCharTok{\%\textgreater{}\%}
    \FunctionTok{mutate}\NormalTok{(}
    \AttributeTok{sentiment\_int =} \FunctionTok{abs}\NormalTok{(vader\_score), }\CommentTok{\# Capture emotional intensity }
    \AttributeTok{likes\_log =} \FunctionTok{log1p}\NormalTok{(like\_count) }\CommentTok{\# Handle skewed distribution}
\NormalTok{  ) }\SpecialCharTok{\%\textgreater{}\%} 
  \CommentTok{\# Plot }
  \FunctionTok{ggplot}\NormalTok{(}\FunctionTok{aes}\NormalTok{(}\AttributeTok{x =}\NormalTok{ sentiment\_int, }\AttributeTok{y =}\NormalTok{ likes\_log)) }\SpecialCharTok{+}
  \FunctionTok{geom\_point}\NormalTok{(}\AttributeTok{alpha =} \FloatTok{0.3}\NormalTok{, }\AttributeTok{color =} \StringTok{"darkblue"}\NormalTok{) }\SpecialCharTok{+}
  \FunctionTok{geom\_smooth}\NormalTok{(}\AttributeTok{method =} \StringTok{"lm"}\NormalTok{, }\AttributeTok{se =} \ConstantTok{FALSE}\NormalTok{, }\AttributeTok{color =} \StringTok{"red2"}\NormalTok{) }\SpecialCharTok{+}
  \FunctionTok{labs}\NormalTok{(}
    \AttributeTok{title =} \StringTok{"Engagement vs Sentiment Intensity"}\NormalTok{,}
    \AttributeTok{x =} \StringTok{"Emotional Intensity"}\NormalTok{,}
    \AttributeTok{y =} \StringTok{"log(1 + likes)"}\NormalTok{) }\SpecialCharTok{+} 
  \FunctionTok{theme\_bw}\NormalTok{()}
\end{Highlighting}
\end{Shaded}

\begin{verbatim}
## `geom_smooth()` using formula = 'y ~ x'
\end{verbatim}

\pandocbounded{\includegraphics[keepaspectratio]{Lab3-hw3_files/Lab3-hw3_files/figure-latex/engagement-1.pdf}}

\textbf{Bonus:} Is there a relationship between sentiment intensity and
engagement? What does that suggest about what gets amplified on Bluesky?

\begin{center}\rule{0.5\linewidth}{0.5pt}\end{center}

\end{document}
